% Internally reviewed construction; exact and LP verification accompany it.
\section{A fixed-budget boundary between scalar and joint precision}
\label{sec:two-target-boundary}
The scalar half-success floor need not survive adding a second target.
The following construction keeps the query budget at two and makes every
coordinate available at unit cost.

\begin{theorem}[Two-target precision with a fixed two-query budget]
\label{thm:two-target-collapse}
For every integer $m\geq3$, let
\[
 F_m=\{(t,u,v,z_1,\ldots,z_m):0\leq u\leq t\leq1,
             \ 0\leq v\leq1-t,\ 0\leq z_j\leq1\}.
\]
All $m+3$ coordinates are exact unit-cost queries. Define
\[
 a_j=j+\tfrac14,\quad b_j=m+1-j+\tfrac14,\quad
 H=\frac{m(2m+3)}4,\quad M=H-m,
\]
and the affine targets
\[
 T_1=Mt+mu+\sum_j a_jz_j,\qquad
 T_2=M(1-t)+mv+\sum_j b_jz_j.
\]
Both initial target intervals are $[0,2H]$. With at most two queries on
every execution, let success mean that both terminal sharp widths are at
most $H$. Then the optimal worst-record precision-success probabilities are
\[
 S^{\star}_{\mathrm{adaptive}}=1,\qquad
 S^{\star}_{\mathrm{randomized\ batch}}=\frac1m.
\]
Each scalar target separately admits a fixed two-query policy with sure
success. The adversarial record is fixed before any policy randomization.
\end{theorem}
\begin{proof}
The polytope $F_m$ is nonempty, compact and convex. Since
$\sum_j a_j=\sum_j b_j=H$ and $M+m=H$, the extrema of $T_1$ are
$0$ and $2H$, attained at $(t,u,z)=(0,0,0)$ and $(1,1,\boldsymbol1)$.
The corresponding choices $(t,v,z)=(1,0,0)$ and
$(0,1,\boldsymbol1)$ give the same extrema for $T_2$.

Query $t$, choose the smallest index $j$ satisfying
$(j-1)/m\leq t\leq j/m$, and query $z_j$. The exact terminal widths are
\begin{equation}
 W_1=H+mt-a_j,\qquad W_2=H+m(1-t)-b_j.
 \label{eq:two-target-boundary-widths}
\end{equation}
Indeed, conditioning on $t$ leaves $u\in[0,t]$ and $v\in[0,1-t]$;
the unqueried $z$ coordinates remain an independent box. The observed
value of $z_j$ contributes only a constant. The chosen branch gives
$\max(W_1,W_2)\leq H-1/4$, proving adaptive success one.

For the randomized-batch lower bound, choose $j\in\{1,\ldots,m\}$ uniformly before any
observations and query $\{t,z_j\}$. Every feasible $t$ belongs to at least
one branch interval, so success probability is at least $1/m$ for every
record. This lower bound holds over the full continuum of $F_m$.

For the upper bound, give equal adverse mass to the $m$ feasible records
\[
 x_i=\bigl((i-\tfrac12)/m,0,0,0,\ldots,0\bigr),
       \qquad i=1,\ldots,m,
\]
and write $t_i=(i-1/2)/m$. The following table lists the exact terminal
widths on these records for every type of distinct two-coordinate batch:
\[
\begin{array}{c|@{\quad}c@{\quad}c}
\text{queried coordinates}&W_1&W_2\\\hline
\{t,z_j\}&H+mt_i-a_j&H+m(1-t_i)-b_j\\
\{t,u\}&H&H+m(1-t_i)\\
\{t,v\}&H+mt_i&H\\
\{u,v\}&H+M&H+M\\
\{u,z_j\}&H+M-a_j&2H-b_j\\
\{v,z_j\}&2H-a_j&H+M-b_j\\
\{z_j,z_k\},\ j\ne k&2H-a_j-a_k&2H-b_j-b_k
\end{array}
\]
For the first row,
\[
 \max(W_1,W_2)=H+|i-j|-\tfrac34,
\]
so the batch succeeds exactly at $x_j$. The next two rows fail because
$0<t_i<1$, and the fourth fails because $M>0$. For the fifth and sixth,
$a_j,b_j\leq m+1/4<H$, so the displayed $2H$ term exceeds $H$.
Finally,
\[
 (a_j+a_k)+(b_j+b_k)=2m+3<2H \qquad(m\geq3).
\]
At least one of $a_j+a_k,b_j+b_k$ is therefore below $H$, so every
two-$z$ batch fails as well. Thus each deterministic batch succeeds on
at most $1/m$ of the adverse mass. The same holds for any batch mixture;
its minimum success over records cannot exceed this adverse average.

Fewer than two distinct queries, repeated queries, and stopping early do
not evade the bound. Extend any such query set to a distinct two-query
set. Additional observations cannot enlarge either sharp interval, so the
original success set is contained in the extended success set. The table
therefore also bounds every permitted shorter execution. This completes
the matching upper and lower bounds.

Finally, the fixed batch $\{t,z_m\}$ makes $W_1\leq H-1/4$ for every
record, and $\{t,z_1\}$ makes $W_2\leq H-1/4$. The difficulty concerns
serving both targets with the same two queries, not either target alone.
\end{proof}

\begin{corollary}[All pairs of three models, threshold margin, and query cost]
\label{cor:all-pairs-collapse}
The construction can be realized as all three paired squared-loss
comparisons among three nonnegative prediction vectors. The exact success
values remain $1$ and $1/m$ at every common normalized threshold
\[
 \tau\in\left[\frac12-\frac1{8H},\,
                    \frac12+\frac1{8H}\right).
\]
Each target separately remains fixed-batch certifiable, and three fixed
queries suffice jointly throughout this interval. Hence the unbounded
ratio concerns success probabilities at a common two-query budget, not
an unbounded reduction in the number of readings.
\end{corollary}
\begin{proof}
Add $T_3=T_1-T_2$ and put
\[
 d_j=a_j-b_j=2j-m-1,\qquad
 P=\sum_{d_j>0}d_j=\lfloor m^2/4\rfloor.
\]
The core $M(2t-1)+m(u-v)$ ranges from $-H$ to $H$, independently of
the $z$ box, whose contribution ranges from $-P$ to $P$. Thus $T_3$ has
initial width $2(H+P)$. After querying $t$ alone its exact width is
$m+2P$. Moreover,
\[
 (H+P)\left(1-\frac1{4H}\right)-(m+2P)
 =M-P-\frac14\left(1+\frac PH\right)>0,
\]
because $M-P\geq m(m-1)/4\geq3/2$ and $P\leq H$.
The third target therefore meets even the smallest displayed normalized
threshold after the first query.

For the first two targets, the displayed interval is equivalent to an
absolute width threshold $q=2H\tau\in[H-1/4,H+1/4)$. The adaptive
branches and uniform batch mixture retain their lower guarantees.
On the adverse records, off-diagonal $\{t,z_j\}$ batches have maximum
width at least $H+1/4$. Every other table row has a larger failing width:
at least $H+1/2$ for $\{t,u\},\{t,v\}$;
$H+M$ for $\{u,v\}$; $H+M-1/4$ for the mixed $u/v,z$ rows;
and $H+M-3/2$ for two-$z$ batches. All exceed $H+1/4$ when $m\geq3$.
Thus the same adverse distribution proves the upper bound $1/m$
throughout the stated half-open interval, also when success includes $T_3$.
The upper endpoint is excluded because an off-diagonal adjacent batch
then succeeds at equality.

The scalar batches from the theorem remain valid throughout the interval,
and $T_3$ needs only $t$. The fixed three-query batch $\{t,z_1,z_m\}$
gives both first widths at most $H-3/2$, while its third width cannot
exceed the width after $t$ alone. Thus three fixed queries suffice;
two fixed queries cannot suffice because their optimal randomized success
is only $1/m<1$.

For an explicit model comparison, write $N=m+3$ and choose the three
nonnegative prediction vectors, in coordinate order $(t,u,v,z_1,\ldots,z_m)$,
\[
 p=(NM/2,0,0,0,\ldots,0),
\]
\[
 q^{(1)}=(NM,Nm/2,0,Na_1/2,\ldots,Na_m/2),\qquad
 q^{(2)}=(0,0,Nm/2,Nb_1/2,\ldots,Nb_m/2).
\]
Then $\operatorname{MSE}(p,x)-\operatorname{MSE}(q^{(r)},x)$ has
the coefficient vector of $T_r$, for $r=1,2$. Additive constants differ
but do not affect widths. The remaining model-pair contrast has the
coefficient vector of $T_1-T_2$, up to orientation. Therefore the guarantee
holds for all three comparisons, each normalized by its own initial width.
\end{proof}

The adaptive normalized margin at half-width is
$1/(8H)=1/[2m(2m+3)]$, which shrinks with $m$. These are exact-query
constructed instances, not a fixed-error robustness result or a new
empirical advantage on the cumulative-meter data. The scalar sparse-cover
theorem concerns one affine target. Here each conditional joint width is
the maximum of two affine functions; its strict failure set can be
disconnected, so the scalar Helly argument does not supply a half-success
floor.
